% Kapitel 6 -- Von dem Thaya Wassermann

\section{Von dem Thaya Wassermann}

\mylettrine{I}{n} dem Thayatal geht seit Menschengedenken eine Sage um, die man sich bei Feuer und Bier erzählt und die man, so ward mir berichtet, mit gesenkter Stimme beendet. Es ist die Sage des \gls{Thaya Wassermann}.

Wer in einer Vollmondnacht sich dem Ufer der Thaya nähert und stillhält, der vernimmt zuerst ein tiefes Rauschen, wie das Atmen eines Wesens, das sich im Schlaf dreht. Dann kommt das Plätschern, leise und regelmäßig, dann das Spritzen. Und in dem Augenblick, da der Mond sein volles Licht auf das Wasser gießt, steigt der Wassermann empor. Er ist grün, so sagen sie: das Haupthaar, der lange Bart, selbst die Haut trägt den Schimmer des Flußgrundes. In seiner Rechten hält er einen goldenen Kamm; mit diesem kämmt er sein Haar, Zug für Zug, bedächtig wie einer, der Zeit hat ohne Ende, und mit jedem Zug fallen silberne Tropfen auf das Ufer, glänzend wie hingestreute Münzen.

Die Tropfen blieben, so heißt es, bis zum ersten Hahnenschrei liegen. Wer sie aufhob, ehe das Tageslicht auf sie fiel, fand wahres Silber in den Händen; mancher Hardegger, so versichern mir die älteren Leute des Ortes, sei auf diese Weise zu Wohlstand gelangt, und die Namen solcher Glücklichen werden noch heute genannt, wenn man die Saga erzählt.

Doch die Sage trägt noch etwas anderes in sich, und das ist der Teil, den die Mütter ihren Kindern einschärften: dass der Wassermann auch lauert. Er liegt manchmal reglos im Schilf, kaum zu unterscheiden von den schlammigen Halmen, und wer zu nahe ans Ufer tritt, den packt er und zieht ihn in die Tiefe, ohne Wort, ohne Vorwarnung. Am liebsten aber, so heißt es, lockt er die Kinder. Ihnen zeigt er durch das spiegelnde Wasser das Bild einer versunkenen Stadt: enge Gassen, bunte Trachten, fremde Menschen, die fröhlich umherwandeln, als hätten sie nie gehört, dass die Welt über ihnen existiert. In der Mitte dieser Stadt hängt eine große silberne Glocke; keinen Laut giebt sie von sich, aber die Kinder bilden sich ein, ihr Klingen zu hören. Da spricht der Wassermann mit sanfter Stimme: \enquote{Wollt ihr nicht sehen, wie es dort unten aussieht?} Und manche Kinder, so heißt es, folgten ihm. Keines kehrte wieder.

Diese Erzählung hörte \gls{Kuonrat von Steinare} von einem alten Stallknecht der Burg \gls{Hardegg}, einem zahnlosen Mann, der sie bei Bier am Feuer erzählte mit dem nachlässigen Ton des Mannes, der meint, sein Zuhörer kenne die Sache schon. Aber \gls{Kuonrat von Steinare} kannte sie nicht.

Er hörte zu, und zwar mit jener gespannten Aufmerksamkeit, die ihn nur zwei Dinge erregten: die Aussicht auf Reichtum und die Aussicht auf das Übernatürliche, das er zugleich fürchtete und begehrte. Diese Erzählung enthielt beides. Silbertropfen am Ufer, niemand sonst zugegen, vor dem ersten Licht aufzuheben: er rechnete nach, wie er immer rechnete, wenn Geld im Spiel war, und das Ergebnis seiner Rechnung gefiel ihm.

Dass der Wassermann Kinder in die Tiefe zog, beunruhigte ihn wenig. Er war kein Kind. Er war ein Knappe, waffengeübt, und hatte überdies, so sagte er sich, keine Absicht, das Silber aus dem Wasser zu fischen; er würde es am Ufer aufheben. Was den Wassermann selbst betraf: er war ein Wesen der Nacht, ein Feind der Ahnungslosen, aber er hatte auch einen goldenen Kamm. Und wer einen goldenen Kamm besaß, den konnte man vielleicht stellen.

In der Nacht nach jenem Abend ließ er den \gls{Knecht} zu sich rufen.

\enquote{Wir gehen heute Nacht zur Thaya}, sprach er zu ihm, \enquote{bei Vollmond. Bring ein starkes Seil.}

Der Knecht sah ihn an. \enquote{Wozu ein Seil, Herr?}

\enquote{Du wirst es sehen.} \gls{Kuonrat von Steinare} wandte sich ab, denn er mochte nicht erklären, was er selbst noch nicht genau ausgeformt hatte. Aber er war entschlossen, und der Entschluss genügte ihm.

Es war eine stille Nacht; der Mond stand rund und gelblichgroß über dem Thayatal und legte sein Licht auf das schwarze Wasser, das lautlos unter der Furt dahinfloß. Die beiden zogen durch das feuchte Gras hinab zum Ufer, der Knabe voran, der Knecht hinter ihm, das Seil über die Schulter gehängt. Das Schilf rauschte leise. Der Fluß roch nach Schlamm und Kälte.

Sie warteten. Eine Stunde verging, vielleicht länger; der Mond wanderte, das Wasser raunte, aber das Plätschern blieb aus und der Wassermann stieg nicht empor.

\enquote{Er kommt nicht}, sagte der \gls{Knecht} nach einer Weile, leise, beinahe flüsternd. \enquote{Vielleicht ist heute nicht die rechte Nacht.}

\enquote{Er kommt}, antwortete \gls{Kuonrat von Steinare}, ohne sich umzusehen. Er stand am Rand des Schilfs und starrte auf das Wasser, die Arme vor der Brust verschränkt, als läge die bloße Kraft seines Willens wie ein Befehl über dem Fluß. Aber der Fluß kümmerte sich nicht um Befehle.

Nach weiterer Zeit sprach er: \enquote{Du wirst ins Wasser steigen. Er kommt nicht, weil hier niemand ist, den er locken kann. Geh hinein und lock ihn heraus.}

Der \gls{Knecht} fuhr zurück. \enquote{Herr, der Wassermann zieht die ins Tiefe, die ins Wasser steigen.}

\enquote{Er zieht Kinder in die Tiefe}, sprach \gls{Kuonrat von Steinare}. \enquote{Kinder, die er täuscht. Du bist kein Kind. Und du läufst nicht weg.} Er dachte einen Augenblick nach. \enquote{Wir binden dir das Seil um den Leib. So kannst du nicht untergehen; ich halte das Ende vom Ufer aus.}

Der \gls{Knecht} sah auf das schwarze Wasser, das im Mondlicht glänzte wie ein Spiegel aus Zinn. Sein Gesicht zeigte das stumme Kämpfen des Unfreien, der einen Befehl nicht ausführen will und doch nicht weiß, wie er Nein sagen soll. \enquote{Herr, wenn die Strömung mich packt, seid Ihr allein nicht stark genug, mich zu halten.}

\enquote{Ich habe gesagt, ich halte das Seil}, sagte \gls{Kuonrat von Steinare} in einem Ton, der keine weiteren Überlegungen zuließ.

So band man dem Knecht das Seil um die Hüften, Knoten auf Knoten, und er stieg ins Wasser. Die Kälte fuhr ihm in die Beine wie ein Stich; er tastete mit den Füßen den Grund, schlüpfrig von Steinen und Schlamm, und trat tiefer, bis das Wasser ihm an die Brust stieg. Der Mond lag auf der Oberfläche und zerschnitt sich in tausend flimmernde Stücke.

Dann verlor er den Boden.

Die Strömung packte ihn mit einem Griff, der ihm den Atem nahm; er wurde seitwärts gerissen, tauchte auf, schrie, trank Wasser, kämpfte und taumelte. Das Seil spannte sich. Am Ufer stand \gls{Kuonrat von Steinare} und schrie Anweisungen ins Dunkel: \enquote{Greif nach ihm! Halt dich aufrecht! Laß ihn nicht entkommen!} Ob er dabei das Seil festhielt oder losgelassen hatte, um besser rufen zu können, das überliefern mir die Berichte nicht.

Was in jener Nacht am Wasser genau geschah, das übersteigt meine sichere Kunde. Die Hardegger aber erzählen es seither so: dass der \gls{Knecht} mit dem Wassermann selbst gerungen habe, Arm in Arm mit dem grünen Geist, und dass \gls{Kuonrat von Steinare} am Ufer gestanden und ihm zugefeuert habe wie ein Feldherr, der seine Männer in die Schlacht schickt und selbst auf der Anhöhe bleibt.

Schließlich, nach langem Ringen, kam der \gls{Knecht} aus dem Wasser. Er war durchnäßt und zerschunden, das Kleid zerrissen, das Haar verklebt mit Schlamm; er zitterte, und seine Füße versagten ihm, als er das Ufer betrat, so dass er in die Knie sank. In seinen Händen aber hielt er, so versichern mir übereinstimmend die Leute, die davon berichten, eine große silberne Glocke. Sie war kalt und schwer und glänzte im Mondlicht wie etwas, das nicht in diese Welt gehört. Kein Ton kam aus ihr; sie schwieg, wie Dinge schweigen, die dem Wasser entstammen und noch nicht wissen, dass sie an die Luft gekommen sind.

\gls{Kuonrat von Steinare} starrte sie an. Er streckte die Hand aus und berührte das Silber, und sein Gesicht zeigte, so sagen jene, die es gesehen haben wollen, etwas, das er selten zeigte: Staunen ohne Hintergedanken, das Staunen des Kindes. Einen Augenblick lang stand er still.

Dann kam die Morgendämmerung. Sie kroch langsam über den Rand der Hügel, erst grau, dann rötlich, und mit dem ersten wirklichen Licht begann die Glocke zu vergehen. Sie zerging nicht wie Wachs, das tropft; sie wurde kleiner und blasser und durchsichtiger, wie ein Bild, das man aus dem Gedächtnis verliert und dem man nachschaut und dabei weiß, dass man es nicht halten kann. Ehe die Sonne über den Horizont stieg, war sie verschwunden. In den Händen des \gls{Knecht}es blieb Leere.

Mancher Hardegger sagt, man habe an jenem Morgen noch einige silberne Tropfen im Schilf gefunden, die im Gras lagen und verschwanden, als man sie aufhob. Ich berichte es nur so, wie es mir erzählt ward, bestätigen kann ich es jedoch nicht.

\gls{Kuonrat von Steinare} aber wandte sich voll Zorn an den Knecht: \enquote{Du bist schuld, dass sie weg ist! Du hast sie nicht festgehalten. Du bist gezittert, und das Silber hat dein Zittern gespürt und ist gebrochen.}

Der \gls{Knecht} schwieg. Er war zu erschöpft, um zu widersprechen, und bereits zu erfahren, um es zu versuchen.

Ob der Knabe in jenem Augenblick wirklich glaubte, was er sagte, das weiß ich nicht. Vielleicht glaubte er es; vielleicht suchte er, wie Menschen zu suchen pflegen, wenn das Glück ausbleibt, einen Schuldigen außerhalb seiner selbst. Der \gls{Knecht} war bei der Hand; er hatte gezittert; und dass er die Glocke überhaupt aus dem Wasser gebracht hatte, das geriet in Vergessenheit.

Zurück auf der Burg \gls{Hardegg} ließ \gls{Kuonrat von Steinare} den \gls{Knecht} in den Schweinestall sperren. Er tat es öffentlich, damit Jedermann sehe, was aus dem wird, der versagt. Die Leute sahen es und schwiegen, wie Leute schweigen, wenn das Urteil ungerecht ist, aber die Gelegenheit fehlt, es zu sagen.

Die Sage des \gls{Thaya Wassermann} aber blieb und blieb erzählt, und wer sie heute hört, findet darin manchmal auch die Geschichte des Knaben, der andere schickte, sein Glück zu holen, und der sich wunderte, dass es nicht hielt.
